% =====================================================================
% LLM Inference Optimization and Scaling to High-Concurrency Serving
% A Six-Pillar Survey of the 2023-2026 Systems Landscape
% Author: Balamurugan Balakreshnan (Microsoft)
%
% arXiv-ready preprint. Compiles with a standard TeX Live install
% (the same toolchain arXiv uses). Build with:
%     pdflatex Bala-Inference-Scaling-Survey-arXiv.tex
%     pdflatex Bala-Inference-Scaling-Survey-arXiv.tex   (run twice for refs/links)
% No .bib file or bibtex pass is required: references are inline below.
% =====================================================================
\documentclass[11pt]{article}

\usepackage[margin=1in]{geometry}
\usepackage{mathptmx}            % Times-like body font
\usepackage[T1]{fontenc}
\usepackage[utf8]{inputenc}
\usepackage{microtype}
\usepackage{booktabs}
\usepackage{tabularx}
\usepackage{ragged2e}
\usepackage{enumitem}
\usepackage{xcolor}
\usepackage{titlesec}
\usepackage{abstract}
\usepackage[hidelinks]{hyperref}

\definecolor{accent}{HTML}{7A1F2B}
\titleformat{\subsection}{\normalfont\bfseries\color{accent}}{\thesubsection}{0.6em}{}

% Tighter caption look
\usepackage{caption}
\captionsetup{font=small,labelfont=bf,labelsep=period}

\newcolumntype{Y}{>{\RaggedRight\arraybackslash}X}

\title{\bfseries LLM Inference Optimization and Scaling to High-Concurrency Serving:\\[2pt]
\large A Six-Pillar Survey of the 2023--2026 Systems Landscape}

\author{Balamurugan Balakreshnan\\
\normalsize Microsoft --- Principal Cloud Solution Architect\\
\normalsize \texttt{babal@microsoft.com}}

\date{Preprint $\cdot$ Survey track $\cdot$ Derived from the author's technical report dated 28 May 2026}

\begin{document}
\maketitle

\begin{abstract}
\noindent
Modern large language model (LLM) inference breaks nearly every assumption baked into classical
request--response serving infrastructure. A single user request is two qualitatively different jobs
glued together: a compute-bound \emph{prefill} that performs dense matrix multiplication over the entire
prompt, and a memory-bandwidth-bound \emph{decode} loop of hundreds of low-arithmetic-intensity steps.
We survey the 2023--2026 systems literature and argue that it has converged on a single architectural
conclusion---stop running prefill and decode on the same GPU---and, more broadly, on six mutually
reinforcing structural moves: (1)~prefill--decode disaggregation, (2)~block-paged KV management with
cross-request prefix reuse, (3)~iteration-level continuous batching with chunked prefill, (4)~low-precision
numerics (FP8/FP4 and weight-only INT4), (5)~sparse computation via Mixture-of-Experts, expert parallelism,
and multi-head latent attention, and (6)~hardware-aware kernel co-design. We synthesize these into a single
taxonomy, show how they compose at million-user scale into a KV-cache-centric distributed system, and connect
the architecture to SLO-driven capacity planning and to production stacks such as Azure AI Foundry, NVIDIA
NIM, TensorRT-LLM, and Dynamo. We close with open problems in inference-aware model--hardware co-design.

\vspace{6pt}
\noindent\textbf{Index terms---}large language models, inference serving, KV cache, PagedAttention,
prefill--decode disaggregation, continuous batching, quantization (FP8/FP4), Mixture-of-Experts,
multi-head latent attention, GPU systems, SLO-driven capacity planning.
\end{abstract}

% =====================================================================
\section{Introduction}
LLM inference is the textbook case of a workload that defeats general-purpose serving infrastructure.
Treating a generation request as an atomic, uniform unit of work---the implicit model behind most web-era
autoscalers and load balancers---leaves enormous amounts of GPU throughput on the table and produces latency
distributions that violate interactive service-level objectives (SLOs). The core reason is structural: the
two phases of a decoder-only transformer request have opposite hardware profiles, opposite scaling laws, and
opposite failure modes.

This survey makes three contributions. First, we consolidate the fast-moving 2023--2026 systems literature
into a \emph{six-pillar taxonomy} that classifies every major optimization by the bottleneck it removes and
the phase it targets. Second, we show how the pillars compose---rather than compete---so that a production
stack adopts all six simultaneously, yielding a KV-cache-centric distributed system rather than a monolithic
model server. Third, we connect the architecture to capacity planning and cost, framing serving design as the
joint optimization of two SLOs (time-to-first-token and time-per-output-token) under a fixed GPU budget.

% =====================================================================
\section{Background: The Two-Phase Anatomy of a Generative Request}
Every decoder-only transformer inference proceeds in two phases. \textbf{Prefill} performs a single forward
pass over all $N$ prompt tokens; it is GEMM-heavy and compute-bound, and it materializes the key--value (KV)
cache of size proportional to $2\,L\,H\,d\,N$ (layers, heads, head-dimension, tokens). \textbf{Decode} then
runs $T$ iterations, generating one token at a time; each step is a GEMV that reads the entire KV cache plus a
new query, making it memory-bandwidth-bound and poorly served by the same hardware configuration that makes
prefill efficient.

Two SLOs govern user-visible latency. \textbf{Time-To-First-Token (TTFT)} is the duration of prefill plus
queue wait; it dominates the perceived responsiveness of chat and the freshness of retrieval-augmented
pipelines. \textbf{Time-Per-Output-Token (TPOT)}, equivalently time-between-tokens (TBT), is the inter-token
interval during decode; it governs streaming smoothness and total completion time. Because the two phases
optimize for different objectives, every pillar below can be read as an attempt to decouple, overlap, or
amortize them.

\begin{table}[t]
\centering
\caption{The six structural pillars of modern LLM inference serving, classified by mechanism, the bottleneck
each removes, and representative systems.}
\label{tab:pillars}
\small
\begin{tabularx}{\textwidth}{@{}p{2.7cm}YYY@{}}
\toprule
\textbf{Pillar} & \textbf{Mechanism} & \textbf{Bottleneck addressed} & \textbf{Representative systems}\\
\midrule
1. Prefill--decode disaggregation & Separate GPU pools per phase; KV cache moved over RDMA & Phase interference; mismatched SLOs & Splitwise, DistServe, Mooncake\\
2. Block-paged KV + prefix reuse & Paged KV blocks; shared-prefix cache hits & KV fragmentation; redundant prefill & vLLM (PagedAttention), SGLang (RadixAttention)\\
3. Continuous batching + chunked prefill & Iteration-level scheduling; prefill piggy-backed on decode & Head-of-line blocking; GPU idle in decode & Orca, Sarathi-Serve\\
4. Low-precision numerics & FP8/FP4 compute; weight-only INT4 & Compute and bandwidth ceiling & FlashAttention-3 FP8, AWQ, GPTQ, SmoothQuant\\
5. Sparse computation & MoE + expert parallelism; latent attention (MLA) & Dense parameter cost; KV-cache size & DeepSeek-V3, Kimi K2\\
6. Hardware-aware kernels & Fused kernels, CUDA Graphs, compiler co-design & Gap to the hardware roofline & FlashAttention-1/2/3, TensorRT-LLM, Dynamo\\
\bottomrule
\end{tabularx}
\end{table}

% =====================================================================
\section{A Six-Pillar Taxonomy of Inference Optimization}

\subsection{Prefill--Decode Disaggregation}
The single most consequential architectural move is to stop co-locating the two phases. Disaggregated serving
runs prefill and decode on separate GPU pools, each sized and scheduled for its own SLO (TTFT for prefill,
TPOT for decode), and transfers the KV cache between them over high-bandwidth RDMA. Splitwise~\cite{splitwise},
DistServe~\cite{distserve}, and Mooncake~\cite{mooncake} each formalize a version of this idea. The payoff is
twofold: prefill bursts no longer stall in-flight decode streams, and each pool can independently exploit
hardware best matched to its arithmetic intensity.

\subsection{Block-Paged KV Management and Cross-Request Prefix Reuse}
Naive contiguous KV allocation wastes 60--80\% of cache memory to internal fragmentation.
PagedAttention, introduced in vLLM~\cite{vllm}, manages the KV cache in fixed-size blocks analogous to
virtual-memory paging, driving waste below a few percent and enabling far larger effective batch sizes.
SGLang's RadixAttention~\cite{sglang} generalizes the idea to \emph{cross-request} reuse: by storing the cache
as a radix tree, queries that share a system prompt or document hit the cache and skip prefill
entirely---now widely deployed as ``prefix caching.'' This pillar attacks both fragmentation and redundant
prefill at once.

\subsection{Iteration-Level Continuous Batching and Chunked Prefill}
Request-level (static) batching forces fast requests to wait for the slowest in their batch.
Orca~\cite{orca} introduced iteration-level ``continuous'' batching, admitting and retiring sequences at every
decode step and eliminating head-of-line blocking. Sarathi-Serve~\cite{sarathi} adds chunked prefill, slicing
long prefills into pieces that are piggy-backed onto decode iterations so the GPU stays compute-bound rather
than stalling on memory-bound decode. Together they keep utilization high without sacrificing tail latency.

\subsection{Low-Precision Numerics: FP8, FP4, and Weight-Only INT4}
Reduced precision doubles or quadruples effective compute and shrinks the bandwidth burden that dominates
decode. FP8 has become the production default on Hopper-class hardware with near-lossless quality; FP4 (NVFP4)
on Blackwell~\cite{nim} pushes throughput further by handling activation dynamic range better than integer
formats. Weight-only schemes---AWQ~\cite{awq} (protecting salient weights), GPTQ~\cite{gptq}
(calibration-based), and SmoothQuant~\cite{smoothquant} (migrating activation outliers into weights)---target
memory-constrained regimes. A recurring lesson is that kernels matter as much as algorithms: identical
quantized weights can run several times faster under better-tuned CUDA kernels.

\begin{table}[t]
\centering
\caption{Representative low-precision formats and their operating points. Figures are indicative of the
2025--2026 landscape and vary by model and calibration.}
\label{tab:quant}
\small
\begin{tabularx}{\textwidth}{@{}lllp{3.2cm}Y@{}}
\toprule
\textbf{Format} & \textbf{Bits} & \textbf{Target HW} & \textbf{Typical quality retention} & \textbf{Best for}\\
\midrule
FP8 & 8 & Hopper+ & $\sim$99\% & Production GPU serving\\
FP4 / NVFP4 & 4 & Blackwell & near-lossless & Max throughput, newest GPUs\\
AWQ & 4 & GPU & $\sim$95\% & Quality-sensitive (coding, creative)\\
GPTQ & 4 & GPU & $\sim$90\% & Max raw throughput\\
INT4 (weight-only) & 4 & General & $\sim$90--95\% & Memory-constrained deployment\\
\bottomrule
\end{tabularx}
\end{table}

\subsection{Sparse Computation: MoE, Expert Parallelism, and Multi-Head Latent Attention}
Mixture-of-Experts (MoE) activates only a subset of parameters per token, amortizing parameter cost; expert
parallelism shards experts across GPUs so total capacity scales without proportional per-token compute.
Multi-head latent attention (MLA), popularized by DeepSeek-V3~\cite{deepseek} and adopted in subsequent
reasoning models, applies low-rank joint compression to keys and values, cutting the KV-cache footprint
several-fold relative to grouped-query attention~\cite{gqa}---directly relieving the bandwidth bottleneck that
dominates decode. Kimi K2~\cite{kimi} and related systems push the same combination at scale.

\subsection{Hardware-Aware Kernel Co-Design}
The final pillar closes the gap to the hardware roofline. FlashAttention 1--3~\cite{flashattn,flashattn2,flashattn3}
recompute attention in tiles to trade memory for compute and exploit asynchrony; CUDA Graphs remove per-kernel
launch overhead; compiler paths (\texttt{torch.compile}, TensorRT-LLM custom kernels~\cite{trtllm}) and
orchestration layers (NVIDIA Dynamo~\cite{dynamo}) co-design the kernel with the memory hierarchy. None of the
higher-level pillars reach their theoretical gains without this low-level co-design.

% =====================================================================
\section{Scaling Out: The KV-Cache-Centric Distributed System}
At million-user scale the serving stack ceases to be a model server and becomes a distributed system organized
around the KV cache. Mooncake (USENIX FAST 2025 best paper)~\cite{mooncake} is the most explicit articulation
of this thesis: a multi-tier KV-cache pool (HBM $\to$ DRAM $\to$ SSD) interconnected by RDMA at up to
$8\times400$~Gbps. The authors report serving on the order of 100 billion tokens per day, with 115\% and 107\%
more requests than prior systems on A800 and H800 clusters respectively. The same pattern, with
Microsoft-specific framing, appears in Splitwise~\cite{splitwise}. The design lesson is that cache placement,
transfer, and eviction---not raw FLOPs---become the primary scaling variables.

% =====================================================================
\section{Capacity Planning and SLO-Driven Cost Modeling}
Because prefill and decode have distinct SLOs and hardware profiles, capacity planning must be done per phase.
Prefill capacity is sized to hold TTFT under burst arrival; decode capacity is sized to hold TPOT under
sustained concurrency. Disaggregation makes this tractable by letting each pool scale on its own signal. The
practical objective is to minimize GPU-hours per million served tokens subject to a TTFT/TPOT SLO pair---an
optimization in which quantization (Pillar~4) and prefix reuse (Pillar~2) move the cost frontier, while
disaggregation (Pillar~1) and continuous batching (Pillar~3) determine how close a deployment runs to that
frontier.

% =====================================================================
\section{Production Integration}
The six pillars are no longer research artifacts. They are landing in production serving stacks---for example
in Azure AI Foundry via NVIDIA NIM~\cite{nim}, TensorRT-LLM~\cite{trtllm}, and Dynamo~\cite{dynamo}
integrations---where disaggregation, paged/prefix-cached KV management, continuous batching, FP8/FP4 numerics,
MoE/MLA model choices, and fused kernels are composed behind a single managed endpoint. The architectural
convergence surveyed here is therefore also an operational convergence: enterprises increasingly inherit all
six pillars by adopting a modern serving platform.

% =====================================================================
\section{Open Problems and Future Directions}
\begin{itemize}[leftmargin=1.4em,itemsep=2pt]
\item \textbf{Inference-aware architectures.} Models are increasingly designed during pre-training to be cheap
to serve (MLA, native multi-token-prediction heads, FP4-friendly weights). Formalizing this model--hardware
co-design---and its accuracy trade-offs---remains open.
\item \textbf{Cache economics across tiers.} Optimal placement, prefetch, and eviction policies for
HBM$\to$DRAM$\to$SSD KV pools under heterogeneous, bursty traffic are not yet well characterized.
\item \textbf{SLO-aware scheduling under disaggregation.} Jointly scheduling prefill and decode pools to hold
both SLOs while maximizing utilization is an under-explored multi-objective control problem.
\item \textbf{Evaluation and safety at serving time.} Continuous, low-overhead evaluation and red-teaming
embedded in the serving path---not only offline---is required for trustworthy autonomous and agentic workloads.
\end{itemize}

% =====================================================================
\section{Conclusion}
The 2023--2026 literature on LLM inference converges on a coherent architecture rather than a grab bag of
tricks. Six pillars---disaggregation, paged/prefix-reused KV, continuous batching with chunked prefill,
low-precision numerics, sparse MoE/MLA computation, and hardware-aware kernels---compose into a
KV-cache-centric distributed system whose design variables are cache movement and SLO-driven capacity, not raw
FLOPs. Framing serving this way turns an open-ended optimization space into a tractable engineering discipline,
and points future work toward inference-aware co-design across the model, the serving stack, and the hardware.

% =====================================================================
\begin{thebibliography}{99}
\small
\bibitem{orca} G.-I. Yu et al., ``Orca: A Distributed Serving System for Transformer-Based Generative Models,'' \emph{OSDI}, 2022.
\bibitem{vllm} W. Kwon et al., ``Efficient Memory Management for Large Language Model Serving with PagedAttention (vLLM),'' \emph{SOSP}, 2023.
\bibitem{mooncake} R. Qin et al., ``Mooncake: A KVCache-Centric Disaggregated Architecture for LLM Serving,'' \emph{USENIX FAST} (Best Paper), 2025.
\bibitem{sglang} L. Zheng et al., ``SGLang: Efficient Execution of Structured Language Model Programs (RadixAttention),'' \emph{NeurIPS}, 2024.
\bibitem{sarathi} A. Agrawal et al., ``Taming Throughput--Latency Tradeoff in LLM Inference with Sarathi-Serve,'' \emph{OSDI}, 2024.
\bibitem{distserve} Y. Zhong et al., ``DistServe: Disaggregating Prefill and Decoding for Goodput-Optimized LLM Serving,'' \emph{OSDI}, 2024.
\bibitem{splitwise} P. Patel et al., ``Splitwise: Efficient Generative LLM Inference Using Phase Splitting,'' \emph{ISCA}, 2024.
\bibitem{flashattn} T. Dao et al., ``FlashAttention: Fast and Memory-Efficient Exact Attention with IO-Awareness,'' \emph{NeurIPS}, 2022.
\bibitem{flashattn2} T. Dao, ``FlashAttention-2: Faster Attention with Better Parallelism and Work Partitioning,'' 2023.
\bibitem{flashattn3} J. Shah et al., ``FlashAttention-3: Fast and Accurate Attention with Asynchrony and Low-Precision,'' 2024.
\bibitem{awq} J. Lin et al., ``AWQ: Activation-Aware Weight Quantization for LLM Compression and Acceleration,'' \emph{MLSys}, 2024.
\bibitem{gptq} E. Frantar et al., ``GPTQ: Accurate Post-Training Quantization for Generative Pre-Trained Transformers,'' \emph{ICLR}, 2023.
\bibitem{smoothquant} G. Xiao et al., ``SmoothQuant: Accurate and Efficient Post-Training Quantization for LLMs,'' \emph{ICML}, 2023.
\bibitem{deepseek} DeepSeek-AI, ``DeepSeek-V3 Technical Report,'' 2024.
\bibitem{kimi} Moonshot AI, ``Kimi K2,'' technical report, 2025.
\bibitem{gqa} J. Ainslie et al., ``GQA: Training Generalized Multi-Query Transformer Models from Multi-Head Checkpoints,'' \emph{EMNLP}, 2023.
\bibitem{specdec} Y. Leviathan et al., ``Fast Inference from Transformers via Speculative Decoding,'' \emph{ICML}, 2023.
\bibitem{trtllm} NVIDIA, ``TensorRT-LLM,'' software documentation, 2024.
\bibitem{dynamo} NVIDIA, ``Dynamo Inference Orchestration,'' technical documentation, 2025.
\bibitem{nim} NVIDIA, ``NIM Microservices and the NVFP4 Format (Blackwell),'' technical documentation, 2025.
\bibitem{report} B. Balakreshnan, ``LLM Inference Optimization and Scaling to High-Concurrency Serving,'' technical report, 28 May 2026.
\end{thebibliography}

\end{document}
